% Fiscal policy and sectoral spillovers --- academic seminar slides
% Prototype talk for de Beauffort & Rannenberg
%   NBB Working Paper No. 493 (June 2026)
%
% Compile from this directory:
%   pdflatex dbr_seminar.tex
%   latexmk -pdf dbr_seminar.tex
%
% Figure: figures/figure3.png (Python + sequence-jacobian prototype IRF).
% Fallback search path: ../output/ if the PNG is regenerated there.
\documentclass[11pt,aspectratio=169]{beamer}

\usetheme{Madrid}
\usecolortheme{default}
\setbeamertemplate{navigation symbols}{}
\setbeamertemplate{itemize items}[circle]
\setbeamerfont{footnote}{size=\tiny}

\usepackage[T1]{fontenc}
\usepackage[utf8]{inputenc}
\usepackage{amsmath,amssymb}
\usepackage{booktabs}
\usepackage{graphicx}
\usepackage{array}

\graphicspath{{figures/}{../output/}}

% Cite the NBB WP in-text; equations/Table 2 follow the public PDF.
\newcommand{\paper}{de Beauffort and Rannenberg (2026)}

\title[Fiscal policy \& sectoral spillovers]{Fiscal Policy and Sectoral Spillovers}
\subtitle{\texorpdfstring{Two-sector open-economy HANK\\
{\normalsize Seminar notes on \paper}}{Open-economy HANK seminar notes}}
\author[Prototype slides]{\texorpdfstring{Open-source notes on NBB Working Paper No.~493\\
{\small not the authors' replication package}}{Open-source prototype slides (not the authors' package)}}
\institute[MacroModellerBot]{\texttt{de-beauffort-rannenberg/} in MacroModellerBot}
\date{September 2026}

\begin{document}

%----------------------------------------------------------------------
\begin{frame}[plain]
  \titlepage
  \vspace{-1.2em}
  \begin{center}
    {\footnotesize
    Source paper: \paper, NBB WP~493 (June 2026),\\
    ``Fiscal policy and sectoral spillovers in open-economy HANK.''}
  \end{center}
\end{frame}

%----------------------------------------------------------------------
\begin{frame}{Outline}
  \tableofcontents
  \vspace{0.8em}
  {\footnotesize Two tracks, in that order: the \emph{paper} (literature, two-sector
  open HANK, claims), then what this repo's Python / sequence-jacobian
  \emph{prototype} actually does.}
\end{frame}

%======================================================================
\section{Literature}
%----------------------------------------------------------------------
\begin{frame}{Why fiscal non-tradables?}
  \begin{block}{Core question \paper}
    Government spending falls disproportionately on \emph{non-tradable services}.
    Empirically, a spending shock still raises private consumption (goods more
    than services), raises goods \emph{and} services output, \emph{lowers} the
    relative price of goods, and \emph{worsens} net exports.
  \end{block}
  \begin{itemize}
    \item Nested two-sector open RANK with the same $G$ bias does \emph{not}
          deliver that package: $C$ falls, and when $G$ is purely NT the
          goods-output co-movement disappears.
    \item The paper asks whether \textbf{household heterogeneity} plus
          \textbf{trade openness} can generate the consumption-led spillover
          from fiscal non-tradables into the goods sector.
    \item Object: \textbf{sectoral fiscal multipliers} in an open economy---not
          a one-sector closed-economy multiplier, and not transfers/FTPL.
  \end{itemize}
\end{frame}

%----------------------------------------------------------------------
\begin{frame}{Sectoral fiscal puzzles: RANK, TANK, open HANK}
  \begin{itemize}
    \item \textbf{Open RANK, sectoral.}
          Monacelli and Perotti (2008, 2010): $C$ falls in both sectors;
          when $G$ is services-intensive, $Y_T$ contracts (negative output
          co-movement). Openness can help relative prices, not that package.
          This paper adds NX and a pre-COVID sample.
    \item \textbf{Fiscal HANK / IKC.}
          High iMPCs reverse RANK crowding-out of $C$ (Gal\'i,
          L\'opez-Salido and Vall\'es 2007 TANK; Auclert, Rognlie and
          Straub 2024 IKC---the household block here). Sticky information
          (Auclert, Rognlie and Straub 2020; $\vartheta=0.935$) for humps;
          signs survive under RE.
    \item \textbf{Open fiscal HANK, one sector.}
          Druedahl et al.\ (2024): import leakage offsets part of the HANK
          multiplier. Aggarwal, Auclert, Rognlie and Straub (2023): high
          iMPCs $\Rightarrow$ persistent CA deficits. Missing piece: T/NT
          relative prices.
    \item \textbf{T / NT structure.}
          Cardi, Restout and Claeys (2020): imperfect labour mobility after
          NT-biased $G$. Corsetti and Dedola (2005): distribution costs in
          retail tradable prices. Composition of $G$ then matters.
  \end{itemize}
  {\scriptsize Contrast, not synthesis: Kaplan--Miyahara (NBER WP 35400) and
  Angeletos, Lian, Wolf and Zhang (NBER WP 35642) are HANK \emph{monetary--fiscal
  / transfer} papers. This talk is NT-biased $G$ and sectoral multipliers.}
\end{frame}

%----------------------------------------------------------------------
\begin{frame}{RANK vs.\ HANK fiscal transmission}
  \begin{itemize}
    \item \textbf{RANK.}
          Ricardian Euler. Deficit-financed $G$ raises expected real rates
          \emph{and} the PV of taxes, so $C$ falls in both sectors
          (Monacelli--Perotti). $Y_T>0$ then rides on the tradable content
          of $G$; with $\phi_{GT}=0$, that co-movement vanishes. NX barely
          moves.
    \item \textbf{HANK.}
          Uninsurable income risk and $a\ge 0$ produce high iMPCs and
          precautionary saving. NT-biased $G$ raises labour income; constrained
          households spend it, and CES demand spills into goods---even if
          $\phi_{GT}=0$.
    \item \textbf{Openness is not a sideshow.}
          Import intensity plus imperfect labour reallocation dampen wage
          pass-through to tradable prices. $P_T/P_{NT}$ can \emph{fall} as
          $Y_T$ rises; NX can deteriorate while expenditure switching still
          supports $Y_T$.
    \item \textbf{Paper's identification of the mechanism.}
          Nested RANK vs.\ HANK, and a $\phi_{GT}=0$ (pure NT) counterfactual:
          HANK keeps $Y_T>0$; RANK does not.
  \end{itemize}
\end{frame}

%----------------------------------------------------------------------
\begin{frame}{Four stylised facts (S-BVAR)}
  Empirical block (WP \S2): Blanchard--Perotti, $G$ ordered first (Cholesky),
  U.S.\ 1954Q1--2019Q4, nine variables, four lags, NIW prior.
  Shock: orthogonal \textbf{1\% of output} increase in $G$.
  \begin{enumerate}
    \item Private consumption \textbf{rises} (crowding-in), with
          $C_T$ \textbf{larger} than $C_{NT}$.
    \item Output rises in \textbf{goods and services} (despite NT-biased $G$).
    \item The relative price of goods \textbf{falls}.
    \item Net exports \textbf{deteriorate}.
  \end{enumerate}
  \vspace{0.35em}
  {\footnotesize
  Ramey-style military narrative is a poor proxy here (more tradable; anticipation).
  The theoretical $G$ path is the empirical IRF (``does not follow a hump-shaped
  path'', \S4.3). This repo uses a 20-quarter path digitised from Figure~3's
  $G$ panel \textbf{[PROVISIONAL]}, not the authors' series.}
\end{frame}

%======================================================================
\section{The model}
%----------------------------------------------------------------------
\begin{frame}{Two-sector open HANK at a glance}
  \begin{itemize}
    \item Domestic production: \textbf{tradables} ($T$, goods) and
          \textbf{non-tradables} ($NT$, services). Linear technology
          $Y_S=N_S$; flexible goods prices, sticky wages only (baseline).
    \item Incomplete-markets households (Auclert--Rognlie--Straub 2024);
          hours set by sectoral unions ($n_{S,t}(i)=N_{S,t}$). Nested RANK:
          same aggregates, Euler instead of the HA Jacobian.
    \item Fiscal $G$ is CES-split with $\phi_{GT}=0.25$ (OECD COFOG tradable
          content); HA-LB / RA-LB set $\phi_{GT}=0$.
    \item Open-economy \emph{block} with foreign variables at SS, calibrated
          to the U.S.\ (12\% import/GDP): CES home/foreign, Corsetti--Dedola
          distribution, UIP with an NFA wedge, tax rule on debt, Taylor on CPI.
    \item Solution: perfect-foresight MIT shocks, sequence-space Jacobians,
          linearised (Auclert, Bard\'oczy, Rognlie and Straub 2021).
  \end{itemize}
\end{frame}

%----------------------------------------------------------------------
\begin{frame}{Households (WP \S3.1)}
  \begin{itemize}
    \item Incomplete markets, $a\ge 0$, idiosyncratic $e_t$ AR(1)
          ($\rho_e=0.968$). Post-tax labour income
          \[
            z_t(i)=(w_T N_T+w_{NT}N_{NT}-T_t)\,e_t(i).
          \]
    \item Utility (4) with $\sigma=\varphi=1$. Households choose $c,a$ only;
          hours are union-set. Taxes scale with $e$, so the $G$ shock does
          not reshape the income \emph{distribution}.
    \item CES demands: $T$ vs.\ $NT$ ($\lambda_t=2.15$, $\phi_t=0.623$,
          including a distribution-margin reclassification);
          home vs.\ foreign with trade elasticity $\lambda_d=1$
          (Cobb--Douglas; 12\% import/GDP via $\phi_d=0.57$).
    \item Sticky expectations, share $\vartheta=0.935$ that does not update
          (paper eq.\ 13 / Auclert--Rognlie--Straub mapping). Used for
          \emph{humps}; qualitative signs are not tied to it.
    \item Nested RANK: same budget, representative Euler, same $r$
          ($\beta_{RA}=0.995$ vs.\ $\beta_{HA}=0.986$).
  \end{itemize}
\end{frame}

%----------------------------------------------------------------------
\begin{frame}{Production, labour, distribution}
  \begin{itemize}
    \item Flexible prices: $P_Z=W_S$ (eq.\ 14). Sticky-price Appendix~C is
          not in the baseline DAG (avoids countercyclical profits).
    \item Imperfect labour mobility: Horvath aggregator (9),
          $\lambda_{\ell}=1$. Sectoral wage Phillips curves (12), slope
          $\kappa_w=0.0044$ (Lind\'e et al.\ 2016), $\theta_w=3$ (50\% markup).
    \item Two reasons $P_T/P_{NT}$ falls: (i)~the foreign component of
          tradable prices; (ii)~NT-biased $G$ plus limited reallocation, so
          NT wages rise more.
    \item Distribution (Corsetti--Dedola, eqs.\ 28--31): local NT services
          enter retail prices of domestic and imported tradables
          ($\mu_d=0.266$, $\mu_f=0.754$). Import NT content $\approx 2\times$
          domestic, which \emph{cushions} NX after appreciation.
    \item Output identities (34)--(35): $Y_T$ is domestic tradables plus
          exports; $Y_{NT}$ is services plus distribution margins.
  \end{itemize}
\end{frame}

%----------------------------------------------------------------------
\begin{frame}{Fiscal, monetary, open economy}
  \begin{align}
    T_t &= T_{ss}+\kappa_b(B_{t-1}-B_{ss})
      \tag{tax, 21} \\[0.2em]
    i_t &= \zeta\, i_{t-1}+(1-\zeta)(i^\ast+\kappa_\pi\pi_t)
      \tag{Taylor, 22}
  \end{align}
  \begin{itemize}
    \item $G$ exogenous; CES split (16)--(19) with $\phi_{GT}=0.25$,
          $\phi^G_d=1$ (no public import leakage in baseline).
          Budget (20). $\kappa_b=0.027$, $\zeta=0.9$, $\kappa_\pi=1.5$.
    \item Mutual fund / UIP (23)--(26). Paper:
          $\Gamma_t=\exp(-10^{-9}\,\mathrm{nfa}_t)$
          (Schmitt-Groh\'e--Uribe NFA closer).
    \item Exports $X_t=C^\ast_{F,t}/(1+\mu_f)$; imports through the
          distribution wedge. Goods/asset clearing (36)--(37).
    \item SSJ DAG: unknowns $(Y_T,Y_{NT},w_T,w_{NT},\pi,Q,i,B,\mathrm{nfa})$;
          $\mathtt{goods\_NT}$ dropped by Walras (household + government +
          NFA laws of motion).
  \end{itemize}
\end{frame}

%----------------------------------------------------------------------
\begin{frame}{Table 2 (paper) vs.\ what we feed the computer}
  \begin{center}
  {\footnotesize
  \begin{tabular}{@{}lrlr@{}}
    \toprule
    Param.\ (Table 2) & Value & Param.\ & Value \\
    \midrule
    $\beta_{RA}$, $\beta_{HA}$ & 0.995, 0.986 & $\phi_t$, $\lambda_t$ & 0.623, 2.15 \\
    $\theta_w$, $\kappa_w$ & 3, 0.0044 & $\phi_d$, $\lambda_d$ & 0.57, 1 \\
    $\sigma$, $\varphi$ & 1, 1 & $\mu_f$, $\mu_d$ & 0.754, 0.266 \\
    $\rho_e$, $\vartheta$ & 0.968, 0.935 & $\zeta$, $\kappa_\pi$ & 0.9, 1.5 \\
    $\phi_{GT}$, $\kappa_b$ & 0.25, 0.027 & $\lambda_{\ell}$, $\phi^G_d$ & 1, 1 \\
    \bottomrule
  \end{tabular}
  }
  \end{center}
  \vspace{0.3em}
  \begin{itemize}
    \item All of the above are \textbf{[VERIFIED]} in \texttt{TABLE2}.
    \item \textbf{[PROVISIONAL]} (NOTES.md; later slide): $\sigma_e=0.35$
          rather than Table~2's $0.58$; $\psi_{\mathrm{nfa}}=10^{-3}$ rather
          than $10^{-9}$; $G/Y$ and $\phi_{\ell}$ implied by adding-up.
    \item Estimation of the S-BVAR itself is \textbf{not} in this repo.
  \end{itemize}
\end{frame}

%======================================================================
\section{Contributions}
%----------------------------------------------------------------------
\begin{frame}{What the NBB WP adds}
  Relative to one-sector HANK and to open-economy RANK fiscal:
  \begin{itemize}
    \item An \textbf{empirical} four-fact package, including NX and a
          $C_T>C_{NT}$ split, for a $G$ shock concentrated in NT services
          (S-BVAR through 2019Q4).
    \item A \textbf{two-sector open HANK}: uninsurable income risk and
          precautionary saving produce a persistent income-driven expansion
          in $C$; import intensity and limited labour reallocation dampen
          wage pass-through to goods prices.
    \item \textbf{Expenditure switching} then yields $Y_T>0$ even as NX
          worsens---and even when $\phi_{GT}=0$. Nested RANK fails that
          test. TANK (App.~B) can raise $C$ on impact but lacks the
          precautionary hump.
    \item Quantitative message: \textbf{household heterogeneity and trade
          openness jointly} shape sectoral fiscal multipliers. Composition
          of $G$ is not a sideshow; more open economies leak more of $Y_T$
          (App.~D).
  \end{itemize}
\end{frame}

%----------------------------------------------------------------------
\begin{frame}{What this open-source prototype adds}
  Honest split vs.\ the published paper:
  \begin{itemize}
    \item \textbf{Runnable Python + sequence-jacobian} two-sector open HANK
          / nested RANK DAG, linear MIT-shock IRFs, and a Figure~3 layout
          (HA, RA, HA-LB, RA-LB, digitised S-BVAR overlay).
    \item Table~2 parameters in \texttt{calibration.py}; qualitative sign
          checks in \texttt{test\_qualitative.py} / \texttt{figure3\_summary.txt}.
    \item Documented \textbf{[PROVISIONAL]} closures needed to complete the
          DAG (next slides): $\sigma_e$, UIP wedge, digitised $G$ / S-BVAR,
          foreign export demand, distribution on $G_D$.
    \item \textbf{Does not} reproduce author IRF numbers, S-BVAR 68\% bands,
          Figure~4 international panel (as the main artefact), TANK App.~B,
          or Appendices~C--D. No official replication package / S-BVAR data
          were found.
  \end{itemize}
\end{frame}

%======================================================================
\section{Results and discussion}
%----------------------------------------------------------------------
\begin{frame}{Prototype IRF: Figure~3 layout}
  \begin{center}
    \includegraphics[width=\textwidth,height=0.68\textheight,keepaspectratio]%
      {figure3.png}
  \end{center}
  {\scriptsize Python + sequence-jacobian 1.0.0 prototype; 1\% of GDP $G$
  shock, empirical-shaped path; $\sigma_e=0.35$, $\psi_{\mathrm{nfa}}=10^{-3}$.
  S-BVAR markers are hand-digitised from WP 493 Figure~3.
  \textbf{[PROVISIONAL]} not the authors' numerical IRFs.}
\end{frame}

%----------------------------------------------------------------------
\begin{frame}{Qualitative co-movements (this code)}
  Signs from \texttt{output/figure3\_summary.txt} (near-term averages),
  matching the paper's main \emph{message}---not author magnitudes:
  \begin{center}
  {\footnotesize
  \begin{tabular}{@{}lcccc@{}}
    \toprule
    Fact & HA & HA-LB & RA & RA-LB \\
    \midrule
    $C>0$ & yes & yes & no (falls) & no \\
    $Y_T>0$ & yes & yes & yes (from $\phi_{GT}$) & \textbf{no} \\
    $Y_{NT}>0$ & yes & yes & yes & yes \\
    $P_T/P_{NT}$ down & yes & yes & yes (smaller) & yes \\
    NX down & yes & yes & no (tiny $+$) & no \\
    $C_T>C_{NT}$ response & yes & yes & yes & yes \\
    \bottomrule
  \end{tabular}
  }
  \end{center}
  \begin{itemize}
    \item Heterogeneity delivers the consumption-led spillover to tradables;
          RANK does not, especially when $G$ is purely non-tradable.
    \item RA still has $Y_T>0$ at $\phi_{GT}=0.25$: that is $G$'s tradable
          content, not crowding-in of $C$.
  \end{itemize}
\end{frame}

%----------------------------------------------------------------------
\begin{frame}{Mechanism, and what is \emph{not} claimed}
  \begin{itemize}
    \item \textbf{HANK mechanism (paper).}
          NT-biased $G$ $\to$ labour income $\to$ high iMPC consumption
          $\to$ CES demand for goods $\to$ $Y_T>0$ even if $\phi_{GT}=0$.
          Distribution + imperfect mobility: $P_T/P_{NT}$ can fall.
          Openness: NX deteriorates while expenditure switching supports
          $Y_T$.
    \item \textbf{RANK contrast.}
          $C$ falls (higher $r$ and future taxes). With $\phi_{GT}=0.25$,
          $Y_T$ can still rise from the tradable content of $G$; with
          $\phi_{GT}=0$, that co-movement disappears.
    \item \textbf{Known shape gaps vs.\ the published figure}
          (NOTES.md; this code):
          HA $C$ is a bit too front-loaded (sticky-info hump milder);
          HA $Y_T$ does not reproduce the empirical hump;
          HA NX stays negative too long (the authors note a medium-run NX
          undershoot);
          RANK NX is a small surplus rather than a small deficit.
  \end{itemize}
\end{frame}

%----------------------------------------------------------------------
\begin{frame}{Why this matters for recent fiscal / open-economy debates}
  \begin{itemize}
    \item \textbf{Composition of $G$.}
          Health, education, administration are services. Multipliers quoted
          on aggregate $G$ hide a T/NT split; $\phi_{GT}=0.25$ is the
          empirically relevant bias, not a knife-edge.
    \item \textbf{Crowding-in vs.\ crowding-out.}
          The S-BVAR sample ends in 2019Q4. The \emph{mechanism}---high
          iMPCs plus NT-biased $G$---is what travels to later packages, not
          a COVID-period estimate. Kaplan--Miyahara / Angeletos et al.\ speak
          to transfers and monetary--fiscal inflation; different object.
    \item \textbf{Openness and leakage.}
          At U.S.\ import intensity, expenditure switching can raise $Y_T$
          even as NX worsens. Appendix~D: a \emph{more} open economy leaks
          more of that goods-sector gain. Closed-economy HANK misses the
          joint restriction.
    \item \textbf{Relative prices.}
          NT-heavy $G$ can \emph{lower} goods prices relative to
          services---a sectoral relative-price fact, not a one-sector
          inflation multiplier.
  \end{itemize}
\end{frame}

%======================================================================
\section{What the repo prototype achieves}
%----------------------------------------------------------------------
\begin{frame}[fragile]{Runnable artefacts}
  \textbf{Runs today} (Python + sequence-jacobian 1.0.0; Table~2 params unless noted):
  \begin{itemize}
    \item Two-sector open HANK (sticky-info HA) and nested RANK.
    \item $\phi_{GT}=0.25$ baseline and $\phi_{GT}=0$ (LB) variants.
    \item Figure~3 PNG + 20-quarter CSVs + qualitative summary.
    \item Smoke test of HA signs, RANK $C<0$, and RA-LB $Y_T\not>0$.
  \end{itemize}
  {\scriptsize
  \begin{verbatim}
cd de-beauffort-rannenberg
python3 src/run_figure3.py
python3 src/run_figure3.py --sigma-e 0.35 --psi-nfa 1e-3
python3 src/run_figure3.py --T 80 --skip-lb   # faster smoke
python3 src/test_qualitative.py
  \end{verbatim}
  }
  Writes \texttt{output/figure3.png}, \texttt{irf\_\{ha,ra,ha\_lb,ra\_lb\}.csv},
  \texttt{figure3\_summary.txt}.

  \textbf{Does not run:} author S-BVAR estimation, exact $G$ vector,
  Appendix~B TANK, Appendix~C sticky prices, Appendix~D 25\% import share.
\end{frame}

%----------------------------------------------------------------------
\begin{frame}{Calibration knobs and caveats}
  \begin{itemize}
    \item \textbf{$\sigma_e$.}
          Table~2 lists $0.58$ (Auclert et al.\ 2024 quarterly conversion).
          As the SSJ Rouwenhorst SD of $\log e$, that undershoots Table~1
          year-1 iMPC $0.507$ and text HtM $49.7\%$ ($\approx 0.24$ /
          $\approx 20\%$). Default \texttt{-{-}sigma-e 0.35} gets
          iMPC$_1\approx 0.53$, HtM $\approx 49.5\%$ given $\beta_{HA}$ and
          $r=1/\beta_{RA}-1$ \textbf{[PROVISIONAL]}---close, not exact.
    \item \textbf{UIP wedge.}
          Paper $\psi_{\mathrm{nfa}}=10^{-9}$; CLI default
          \texttt{-{-}psi-nfa 1e-3} for a stable linear solve
          \textbf{[PROVISIONAL]}.
    \item \textbf{S-BVAR overlay} in \texttt{empirical\_sbvar.py}: hand-read
          from Figure~3, not the authors' series \textbf{[PROVISIONAL]}.
    \item \textbf{Other closures} (NOTES.md): $G/Y\approx 0.214$ implied by
          12\% import/GDP; $\phi_{\ell}$ set to sectoral hours shares so both
          wage PCs hold at $w=1$; distribution margin on $G_D$ for labour
          income $=$ GDP; SOE export demand with $C^\ast$ fixed; SSJ
          ex-post $r=$ lagged ex-ante.
  \end{itemize}
\end{frame}

%======================================================================
\section{Simplifying assumptions}
%----------------------------------------------------------------------
\begin{frame}{What is on, what is shut off}
  \textbf{In the paper's baseline:}
  \begin{itemize}
    \item Two domestic sectors only; linear technology; flexible prices /
          sticky wages (not sticky prices).
    \item Hours set by unions; incomplete markets with $a\ge 0$; sticky
          expectations on the HA Jacobian (signs not tied to it).
    \item Foreign variables at SS; Cobb--Douglas trade ($\lambda_d=1$);
          Horvath $\lambda_{\ell}=1$; $\phi_{GT}=0.25$; $\phi^G_d=1$.
  \end{itemize}
  \textbf{Extra prototype approximations (not author codes):}
  \begin{itemize}
    \item $\sigma_e=0.35$ vs.\ Table~2 $0.58$; $\psi_{\mathrm{nfa}}=10^{-3}$
          vs.\ $10^{-9}$; digitised $G$ and S-BVAR; provisional foreign
          CES / distribution on $G_D$; SSJ $r$ timing.
  \end{itemize}
  \textbf{Shut off in both (baseline script):}
  Appendix~B TANK; Appendix~C sticky prices; Appendix~D 25\% import share;
  Figure~4 international panel as the main artefact; official S-BVAR data.
\end{frame}

%======================================================================
\section{Extensions}
%----------------------------------------------------------------------
\begin{frame}{Future developments that fit this codebase}
  \begin{enumerate}
    \item \textbf{N-sector sketch} (\texttt{EXTENDING.md} / \texttt{sectors.py}).
          Today \texttt{DOMESTIC\_SECTORS = (T, NT)}; blocks are still
          explicit two-sector objects. Generate per-sector wage PC / hours /
          demand from that tuple instead of copy-pasting $T$/$NT$.
    \item \textbf{Nesting, not a monolith.}
          A third sector $S$ needs an explicit CES / distribution / tradable
          status. Keep trade, UIP, and the HA Jacobian as separate SSJ
          blocks. Walras: still drop \emph{exactly one} sectoral goods
          residual.
    \item \textbf{Paper appendices} (DAG is extensible).
          TANK (App.~B); sticky prices (App.~C); 25\% import share (App.~D);
          Figure~4 international panel as a first-class artefact.
    \item \textbf{Blocked on data/codes.}
          Official replication files, S-BVAR data, exact $G$ IRF vector,
          fully specified foreign block.
    \item \textbf{Later, not this folder.}
          HA-OLG household block: Bard\'oczy--Shanker--Vel\'asquez-Giraldo
          (FEDS 2026-046) is the SSJ cookbook---unused here (infinite horizon).
          Synthesis with the HENK companion: two-sector open HANK fiscal
          plus heterogeneous-expectation NK.
  \end{enumerate}
\end{frame}

%----------------------------------------------------------------------
\begin{frame}{Takeaways}
  \begin{itemize}
    \item NT-biased government spending can still raise $C$, raise $Y_T$
          \emph{and} $Y_{NT}$, lower $P_T/P_{NT}$, and worsen NX---if
          households are incomplete-markets and the economy is open.
    \item Nested RANK fails that package, especially when $\phi_{GT}=0$.
          Heterogeneity, not the tradable content of $G$, is the spillover.
    \item This repo is a transparent Figure~3 starter: Python + SSJ DAG,
          qualitative HANK vs.\ RANK signs, documented provisional
          $\sigma_e$ / UIP / S-BVAR closures---not a numerical replica of
          the authors' IRFs.
    \item The natural extension is a labelled multi-sector open DAG, not a
          one-file N-sector solver.
  \end{itemize}
  \vfill
  {\scriptsize
  \textbf{Selected references.}
  de Beauffort and Rannenberg (2026), NBB WP~493;
  Monacelli and Perotti (2008, 2010);
  Auclert, Rognlie and Straub (2020, 2024);
  Auclert, Bard\'oczy, Rognlie and Straub (2021);
  Druedahl et al.\ (2024); Aggarwal et al.\ (2023);
  Cardi, Restout and Claeys (2020); Corsetti and Dedola (2005);
  Gal\'i, L\'opez-Salido and Vall\'es (2007).
  Compile: \texttt{latexmk -pdf dbr\_seminar.tex} from
  \texttt{de-beauffort-rannenberg/beamer/}.}
\end{frame}

\end{document}
